\chapter{2011年福建卷（理）}

\section{单选题}

\begin{problem}
\(\i\) 是虚数单位，若集合 \(S = \{-1, 0, 1\}\)，则（\quad）

\choices
    {\(\i \in S\)}
    {\(\i^2\in S\)}
    {\(\i^3\in S\)}
    {\(\frac{2}{\i} \in S\)}
\end{problem}

\begin{answer}
B
\end{answer}

\begin{solution}
由虚数单位的定义，\(\i^2=-1\)，所以 \(\i^2\in S\). 又有 \(\i\notin S\)，\(\i^3=-\i\notin S\)，且 \(\frac{2}{\i}=-2\i\notin S\). 因此选 B.
\end{solution}

\begin{problem}
若 \(a \in \R\)，则 “\(a = 2\)” 是 “\((a - 1)(a - 2) = 0\)” 的（\quad）

\choices
    {充分而不必要条件}
    {必要而不充分条件}
    {充要条件}
    {既不充分也不必要条件}
\end{problem}

\begin{answer}
A
\end{answer}

\begin{solution}
当 \(a=2\) 时，\((a-1)(a-2)=0\)，所以“\(a=2\)”是“\((a-1)(a-2)=0\)”的充分条件. 反过来，由 \((a-1)(a-2)=0\) 可得 \(a=1\) 或 \(a=2\)，不能推出 \(a=2\)，所以它不是必要条件. 因此选 A.
\end{solution}

\begin{problem}
若 \(\tan \alpha = 3\)，则 \(\frac{\sin 2\alpha}{\cos^2 \alpha}\) 的值等于（\quad）

\choices
    {2}
    {3}
    {4}
    {6}
\end{problem}

\begin{answer}
D
\end{answer}

\begin{solution}
因为 \(\tan\alpha=3\)，所以 \(\cos\alpha\ne0\). 利用二倍角公式，得
\[
\frac{\sin2\alpha}{\cos^2\alpha}=\frac{2\sin\alpha\cos\alpha}{\cos^2\alpha}=2\tan\alpha=6.
\]
因此选 D.
\end{solution}

\begin{problem}
如图，矩形 \(ABCD\) 中，点 \(E\) 为边 \(CD\) 的中点，若在矩形 \(ABCD\) 内部随机取一个点 \(Q\)，则点 \(Q\) 取自 \(\triangle ABE\) 内部的概率等于（\quad）

\bitmapfigure[max width=0.34\linewidth]{img/2011/fujian_science/q4_fig1.png}

\choices
    {\( \frac{1}{4} \)}
    {\( \frac{1}{3} \)}
    {\( \frac{1}{2} \)}
    {\( \frac{2}{3} \)}
\end{problem}

\begin{answer}
C
\end{answer}

\begin{solution}
设矩形 \(ABCD\) 的面积为 \(S\)，则 \(E\) 到边 \(AB\) 的距离等于矩形的高. 由于 \(E\) 是 \(CD\) 的中点，\(AB=CD\)，所以
\[
S_{\triangle ABE}=\frac{1}{2}AB\cdot h=\frac{1}{2}S.
\]
随机点在矩形内均匀分布，故所求概率为
\[
\frac{S_{\triangle ABE}}{S}=\frac{1}{2}.
\]
因此选 C.
\end{solution}

\begin{problem}
\(\int_{0}^{1}\left(\e^{x} + 2x\right)\mathrm{d}x\) 等于（\quad）

\choices
    {1}
    {\e - 1}
    {\e}
    {\e + 1}
\end{problem}

\begin{answer}
C
\end{answer}

\begin{solution}
\[
\int_0^1\left(\e^x+2x\right)\mathrm{d}x=\left(\e^x+x^2\right)\bigg|_0^1=(\e+1)-(1+0)=\e.
\]
因此选 C.
\end{solution}

\begin{problem}
\((1 + 2x)^{5}\) 的展开式中，\(x^{2}\) 的系数等于（\quad）

\choices
    {80}
    {40}
    {20}
    {10}
\end{problem}

\begin{answer}
B
\end{answer}

\begin{solution}
由二项式定理，\((1+2x)^5\) 中 \(x^2\) 的项来自取两个 \(2x\) 和三个 \(1\)，其系数为
\[
\mathrm C_5^2\cdot 2^2=10\cdot4=40.
\]
因此选 B.
\end{solution}

\begin{problem}
设圆锥曲线 \(T\) 的两个焦点分别为 \(F_{1}\)，\(F_{2}\). 若曲线 \(T\) 上存在点 \(P\) 满足 \(|PF_{1}|: |F_{1}F_{2}|: |PF_{2}| = 4:3:2\)，则曲线 \(T\) 的离心率等于（\quad）

\choices
    {\(\frac{1}{2}\) 或 \(\frac{3}{2}\)}
    {\(\frac{2}{3}\) 或 2}
    {\(\frac{1}{2}\) 或 2}
    {\(\frac{2}{3}\) 或 \(\frac{3}{2}\)}
\end{problem}

\begin{answer}
A
\end{answer}

\begin{solution}
设存在的点 \(P\) 到两焦点的距离分别为 \(4k\) 和 \(2k\)，其中 \(k>0\). 由题意，\(F_1F_2=3k\)，记焦半距为 \(c\)，则 \(2c=3k\)，即 \(c=\frac{3}{2}k\).

若圆锥曲线为椭圆，则两焦点距离之和为 \(2a\)，所以 \(2a=4k+2k=6k\)，从而 \(a=3k\)，离心率为 \(e=\frac{c}{a}=\frac{1}{2}\).

若圆锥曲线为双曲线，则两焦点距离之差为 \(2a\)，所以 \(2a=4k-2k=2k\)，从而 \(a=k\)，离心率为 \(e=\frac{c}{a}=\frac{3}{2}\). 因此选 A.
\end{solution}

\begin{problem}
已知 \(O\) 是坐标原点，点 \(A(-1,1)\)，若点 \(M(x,y)\) 为平面区域 \(\begin{cases}
x+y\geqslant 2,\\
x\leqslant 1,\\
y\leqslant 2
\end{cases}\) 上的一个动点，则 \(\overrightarrow{OA}\cdot\overrightarrow{OM}\) 的取值范围是（\quad）

\choices
    {\([-1,0]\)}
    {\([0,1]\)}
    {\([0,2]\)}
    {\([-1,2]\)}
\end{problem}

\begin{answer}
C
\end{answer}

\begin{solution}
由 \(x+y\geqslant2\) 和 \(y\leqslant2\)，可得 \(x\geqslant0\). 结合 \(x\leqslant1\)，有 \(0\leqslant x\leqslant1\)，且对每个这样的 \(x\)，满足 \(2-x\leqslant y\leqslant2\). 又
\[
\overrightarrow{OA}\cdot\overrightarrow{OM}=(-1,1)\cdot(x,y)=y-x.
\]
在该区域内，\(y-x\) 的最小值在 \((1,1)\) 处取得，为 \(0\)；最大值在 \((0,2)\) 处取得，为 \(2\). 因此取值范围为 \([0,2]\)，选 C.
\end{solution}

\begin{problem}
对于函数 \(f(x)=a\sin x+bx+c\)（其中 \(a\)，\(b\in\R\)，\(c\in\Z\)），选取 \(a\)，\(b\)，\(c\) 的一组值计算 \(f(1)\) 和 \(f(-1)\)，所得出的正确结果一定不可能是（\quad）

\choices
    {4 和 6}
    {3 和 1}
    {2 和 4}
    {1 和 2}
\end{problem}

\begin{answer}
D
\end{answer}

\begin{solution}
令 \(u=a\sin1+b\)，则
\[
f(1)=u+c,\qquad f(-1)=-u+c.
\]
因此两个结果之和为 \(f(1)+f(-1)=2c\)，必为偶整数. 选项 A、B、C 中的两数之和分别为 \(10\)、\(4\)、\(6\)，均为偶整数，而且分别取 \((a,b,c)=(0,-1,5)\)、\((0,1,2)\)、\((0,-1,3)\) 时可以得到对应结果. 选项 D 中 \(1+2=3\) 不是偶整数，不可能出现. 因此选 D.
\end{solution}

\begin{problem}
已知函数 \(f(x)=\e^{x}+x\). 对于曲线 \(y=f(x)\) 上横坐标成等差数列的三个点 \(A\)，\(B\)，\(C\)，给出以下判断：

\begin{circlelist}
\item \(\triangle ABC\) 一定是钝角三角形；
\item \(\triangle ABC\) 可能是直角三角形；
\item \(\triangle ABC\) 可能是等腰三角形；
\item \(\triangle ABC\) 不可能是等腰三角形.
\end{circlelist}

其中，正确的判断是（\quad）

\choices
    {\circled{1}\circled{3}}
    {\circled{1}\circled{4}}
    {\circled{2}\circled{3}}
    {\circled{2}\circled{4}}
\end{problem}

\begin{answer}
B
\end{answer}

\begin{solution}
设三个点的横坐标依次为 \(t-d\)，\(t\)，\(t+d\)，其中 \(d>0\). 记
\[
\Delta_1=f(t)-f(t-d)=\e^{t-d}(\e^d-1)+d,
\]
\[
\Delta_2=f(t+d)-f(t)=\e^t(\e^d-1)+d.
\]
因为
\[
\Delta_2-\Delta_1=\e^{t-d}(\e^d-1)^2>0,
\]
所以 \(0<\Delta_1<\Delta_2\). 在点 \(B\) 处，两条边向量可写为 \((-d,-\Delta_1)\) 和 \((d,\Delta_2)\)，其内积为
\[
-d^2-\Delta_1\Delta_2<0.
\]
故 \(\angle ABC\) 一定为钝角，判断\(\circled{1}\)正确，判断\(\circled{2}\)错误.

三边长满足
\[
AB^2=d^2+\Delta_1^2<d^2+\Delta_2^2=BC^2,
\]
且
\[
AC^2=4d^2+(\Delta_1+\Delta_2)^2>BC^2.
\]
所以三边两两不相等，判断\(\circled{3}\)错误，判断\(\circled{4}\)正确. 因此选 B.
\end{solution}

\section{填空题}

\begin{problem}
运行如图所示的程序，输出的结果是\fillinblank{}.

\bitmapfigure[max width=0.22\linewidth]{img/2011/fujian_science/q11_fig1.png}
\end{problem}

\begin{answer}
3
\end{answer}

\begin{solution}
程序开始时 \(a=1\)，\(b=2\). 执行 \(a=a+b\) 后，\(a=1+2=3\)，随后输出 \(a\)，所以输出结果为 \(3\).
\end{solution}

\begin{problem}
三棱锥 \(P - ABC\) 中，\(PA \perp\) 底面 \(ABC\)，\(PA = 3\)，底面 \(ABC\) 是边长为 \(2\) 的正三角形，则三棱锥 \(P - ABC\) 的体积等于\fillinblank{}.
\end{problem}

\begin{answer}
\(\sqrt{3}\)
\end{answer}

\begin{solution}
底面 \(ABC\) 是边长为 \(2\) 的正三角形，所以
\[
S_{\triangle ABC}=\frac{1}{2}\cdot2\cdot\sqrt{2^2-1^2}=\sqrt{3}.
\]
又因为 \(PA\perp\) 底面 \(ABC\)，故三棱锥的高为 \(PA=3\). 因此体积为
\[
V=\frac{1}{3}S_{\triangle ABC}\cdot PA=\frac{1}{3}\cdot\sqrt{3}\cdot3=\sqrt{3}.
\]
\end{solution}

\begin{problem}
盒中装有形状，大小完全相同的 \(5\) 个球，其中红色球 \(3\) 个，黄色球 \(2\) 个. 若从中随机取出 \(2\) 个球，则所取出的 \(2\) 个球颜色不同的概率等于\fillinblank{}.
\end{problem}

\begin{answer}
\(\frac{3}{5}\)
\end{answer}

\begin{solution}
从 \(5\) 个球中取出 \(2\) 个球的取法总数为 \(\mathrm C_5^2=10\). 取到一个红球和一个黄球的取法数为 \(\mathrm C_3^1\mathrm C_2^1=6\)，所以所求概率为
\[
\frac{\mathrm C_3^1\mathrm C_2^1}{\mathrm C_5^2}=\frac{6}{10}=\frac{3}{5}.
\]
\end{solution}

\begin{problem}
如图，\(\triangle ABC\) 中，\(AB = AC = 2\)，\(BC = 2\sqrt{3}\)，点 \(D\) 在 \(BC\) 边上，\(\angle ADC = 45^\circ\)，则 \(AD\) 的长度等于\fillinblank{}.

\bitmapfigure[max width=0.28\linewidth]{img/2011/fujian_science/q14_fig1.png}
\end{problem}

\begin{answer}
\(\sqrt{2}\)
\end{answer}

\begin{solution}
取直角坐标系，使 \(B=(-\sqrt{3},0)\)，\(C=(\sqrt{3},0)\). 因为 \(AB=AC=2\)，所以 \(A\) 在 \(BC\) 的垂直平分线上，且
\[
AB^2-\left(\frac{BC}{2}\right)^2=2^2-(\sqrt{3})^2=1,
\]
从而可取 \(A=(0,1)\). 设 \(D=(t,0)\). 由于 \(\angle ADC=45^\circ\)，且 \(\overrightarrow{DC}\) 沿 \(x\) 轴正方向，向量 \(\overrightarrow{DA}=(-t,1)\) 与 \(x\) 轴正方向的夹角为 \(45^\circ\)，故 \(-t=1\)，即 \(D=(-1,0)\). 因此
\[
AD=\sqrt{(0+1)^2+(1-0)^2}=\sqrt{2}.
\]
\end{solution}

\begin{problem}
设 \(V\) 是全体平面向量构成的集合. 若映射 \(f:V\to\R\) 满足：对任意向量 \(a=(x_1,y_1)\in V\)，\(b=(x_2,y_2)\in V\)，以及任意 \(\lambda\in\R\)，均有 \(f\left(\lambda a+(1-\lambda)b\right)=\lambda f(a)+(1-\lambda)f(b)\)，则称映射 \(f\) 具有性质 \(P\). 现给出如下映射：

\begin{circlelist}
\item \(f_1:V\to\R\)，\(f_1(m)=x-y\)，\(m=(x,y)\in V\)；
\item \(f_2:V\to\R\)，\(f_2(m)=x^2+y\)，\(m=(x,y)\in V\)；
\item \(f_3:V\to\R\)，\(f_3(m)=x+y+1\)，\(m=(x,y)\in V\).
\end{circlelist}

其中，具有性质 \(P\) 的映射的序号为\fillinblank{}（写出所有具有性质 \(P\) 的映射的序号）.
\end{problem}

\begin{answer}
\(\circled{1}\circled{3}\)
\end{answer}

\begin{solution}
对映射 \(f_1\)，有
\[
f_1\left(\lambda a+(1-\lambda)b\right)=\lambda(x_1-y_1)+(1-\lambda)(x_2-y_2)=\lambda f_1(a)+(1-\lambda)f_1(b).
\]
所以 \(f_1\) 具有性质 \(P\).

对映射 \(f_2\)，取 \(a=(1,0)\)，\(b=(0,0)\)，\(\lambda=\frac{1}{2}\)，则
\[
f_2\left(\frac{1}{2}a+\frac{1}{2}b\right)=f_2\left(\frac{1}{2},0\right)=\frac{1}{4},
\]
而
\[
\frac{1}{2}f_2(a)+\frac{1}{2}f_2(b)=\frac{1}{2}\cdot1+\frac{1}{2}\cdot0=\frac{1}{2}.
\]
两者不相等，所以 \(f_2\) 不具有性质 \(P\).

对映射 \(f_3\)，有
\[
f_3\left(\lambda a+(1-\lambda)b\right)=\lambda(x_1+y_1)+(1-\lambda)(x_2+y_2)+1=\lambda f_3(a)+(1-\lambda)f_3(b).
\]
所以 \(f_3\) 具有性质 \(P\). 因此具有性质 \(P\) 的映射序号为 \(\circled{1}\circled{3}\).
\end{solution}

\section{解答题}

\begin{problem}
已知等比数列 \(\{a_n\}\) 的公比 \(q=3\)，前 \(3\) 项和 \(S_3=\frac{13}{3}\).

\begin{enumerate}
\item 求数列 \(\{a_n\}\) 的通项公式；
\item 若函数 \(f(x)=A\sin\left(2x+\varphi\right)\)（\(A>0\)，\(0<\varphi<\pi\)）在 \(x=\frac{\pi}{6}\) 处取得最大值，且最大值为 \(a_3\)，求函数 \(f(x)\) 的解析式.
\end{enumerate}
\end{problem}

\begin{solution}
\begin{enumerate}
\item 由等比数列前 \(3\) 项和公式，得
\[
S_3=a_1\left(1+q+q^2\right)=a_1(1+3+9)=13a_1=\frac{13}{3},
\]
所以 \(a_1=\frac{1}{3}\). 因此数列的通项公式为
\[
a_n=a_1q^{n-1}=\frac{1}{3}\cdot3^{n-1}=3^{n-2}.
\]

\item 由（1）得 \(a_3=3\). 函数 \(f(x)=A\sin(2x+\varphi)\) 的最大值为 \(A\)，题设给出 \(A=a_3=3\). 在 \(x=\frac{\pi}{6}\) 处取得最大值，故
\[
2\cdot\frac{\pi}{6}+\varphi=\frac{\pi}{2}+2k\pi.
\]
其中 \(k\in\Z\). 由于 \(0<\varphi<\pi\)，左端属于 \(\left(\frac{\pi}{3},\frac{4\pi}{3}\right)\)，所以只能取 \(k=0\)，从而 \(\varphi=\frac{\pi}{6}\). 故
\[
f(x)=3\sin\left(2x+\frac{\pi}{6}\right).
\]
\end{enumerate}
\end{solution}

\begin{problem}
已知直线 \(l:y=x+m\)，\(m\in\R\).

\begin{enumerate}
\item 若以点 \(M(2,0)\) 为圆心的圆与直线 \(l\) 相切于点 \(P\)，且点 \(P\) 在 \(y\) 轴上，求该圆的方程；
\item 若直线 \(l\) 关于 \(x\) 轴对称的直线为 \(l'\)，问直线 \(l'\) 与抛物线 \(C:x^2=4y\) 是否相切？说明理由.
\end{enumerate}
\end{problem}

\begin{solution}
\begin{enumerate}
\item 因为切点 \(P\) 在 \(y\) 轴上，且 \(P\in l\)，所以令 \(x=0\)，得 \(P=(0,m)\). 直线 \(l\) 的斜率为 \(1\)，圆在切点处的半径 \(MP\) 应与 \(l\) 垂直，因此 \(MP\) 的斜率为 \(-1\). 另一方面，
\[
\frac{m-0}{0-2}=-\frac{m}{2}.
\]
所以 \(-\frac{m}{2}=-1\)，即 \(m=2\)，切点为 \(P=(0,2)\). 半径平方为
\[
MP^2=(0-2)^2+(2-0)^2=8.
\]
故圆的方程为
\[
(x-2)^2+y^2=8.
\]

\item 由（1）知 \(l:y=x+2\). 关于 \(x\) 轴对称后，直线 \(l'\) 的方程为 \(y=-x-2\). 将其代入抛物线 \(x^2=4y\)，得
\[
x^2=4(-x-2)
\]
即
\[
x^2+4x+8=0.
\]
该方程的判别式为 \(\Delta=4^2-4\cdot1\cdot8=-16<0\)，所以 \(l'\) 与抛物线没有公共点，因而不相切.
\end{enumerate}
\end{solution}

\begin{problem}
某商场销售某种商品的经验表明，该商品每日的销售量 \(y\)（单位：千克）与销售价格 \(x\)（单位：元/千克）满足关系式 \(y=\frac{a}{x-3}+10(x-6)^2\)，其中 \(3<x<6\)，\(a\) 为常数. 已知销售价格为 \(5\text{ 元/千克}\) 时，每日可售出该商品 \(11\text{ 千克}\).

\begin{enumerate}
\item 求 \(a\) 的值；
\item 若该商品的成本为 \(3\text{ 元/千克}\)，试确定销售价格 \(x\) 的值，使商场每日销售该商品所获得的利润最大.
\end{enumerate}
\end{problem}

\begin{solution}
\begin{enumerate}
\item 将 \(x=5\)，\(y=11\) 代入题设关系式，得
\[
11=\frac{a}{5-3}+10(5-6)^2=\frac{a}{2}+10
\]
所以 \(a=2\).

\item 每日销售该商品获得的利润为销售量乘以每千克的利润，因此
\[
L(x)=(x-3)y=(x-3)\left[\frac{2}{x-3}+10(x-6)^2\right]
 =2+10(x-3)(x-6)^2.
\]
令 \(g(x)=(x-3)(x-6)^2\)，则
\[
g'(x)=(x-6)^2+2(x-3)(x-6)=3(x-6)(x-4).
\]
在 \(3<x<4\) 时，\(g'(x)>0\)；在 \(4<x<6\) 时，\(g'(x)<0\). 所以 \(g(x)\)，从而 \(L(x)\)，先增后减，在 \(x=4\) 时取得最大值. 故销售价格应确定为 \(x=4\) 元/千克.
\end{enumerate}
\end{solution}

\begin{problem}
某产品按行业生产标准分成 \(8\) 个等级，等级系数 \(X\) 依次为 \(1\)，\(2\)，\(\dots\)，\(8\)，其中 \(X\geqslant 5\) 为标准 \(A\)，\(X\geqslant 3\) 为标准 \(B\). 已知甲厂执行标准 \(A\) 生产该产品，产品的零售价为 \(6\text{ 元/件}\)；乙厂执行标准 \(B\) 生产该产品，产品的零售价为 \(4\text{ 元/件}\)，假定甲，乙两厂的产品都符合相应的执行标准.

\begin{enumerate}
\item 已知甲厂产品的等级系数 \(X_1\) 的概率分布如下表所示：

\begin{center}
\begin{tabular}{|c|c|c|c|c|}
\hline
\(X_1\) & 5 & 6 & 7 & 8 \\
\hline
\(P\) & 0.4 & \(a\) & \(b\) & 0.1 \\
\hline
\end{tabular}
\end{center}

且 \(E(X_1)=6\). 求 \(a\)，\(b\) 的值；
\item 为分析乙厂产品的等级系数 \(X_2\)，从该厂生产的产品中随机抽取 \(30\) 件，相应的等级系数组成一个样本. 数据如下：

\begin{center}
\begin{tabular}{cccccccccc}
3 & 5 & 3 & 3 & 8 & 5 & 5 & 6 & 3 & 4 \\
6 & 3 & 4 & 7 & 5 & 3 & 4 & 8 & 5 & 3 \\
8 & 3 & 4 & 3 & 4 & 4 & 7 & 5 & 6 & 7
\end{tabular}
\end{center}

用这个样本的频率分布估计总体分布，将频率视为概率，求等级系数 \(X_2\) 的数学期望；
\item 在前两问的条件下，若以“性价比”为判断标准，则哪个工厂的产品更具可购买性？说明理由.
\end{enumerate}

注：

\begin{circlelist}
\item 产品的“性价比”\(=\frac{\text{产品的等级系数的数学期望}}{\text{产品的零售价}}\)；
\item “性价比”大的产品更具可购买性.
\end{circlelist}
\end{problem}

\begin{solution}
\begin{enumerate}
\item 由概率总和为 \(1\)，得
\[
0.4+a+b+0.1=1
\]
即 \(a+b=0.5\). 又由 \(E(X_1)=6\)，得
\[
5\cdot0.4+6a+7b+8\cdot0.1=6
\]
即 \(6a+7b=3.2\). 联立两式，得 \(a=0.3\)，\(b=0.2\).

\item 样本中等级系数 \(3\)，\(4\)，\(5\)，\(6\)，\(7\)，\(8\) 分别出现 \(9\)，\(6\)，\(6\)，\(3\)，\(3\)，\(3\) 次. 因此将频率视为概率时，
\[
E(X_2)=3\cdot\frac{9}{30}+4\cdot\frac{6}{30}+5\cdot\frac{6}{30}+6\cdot\frac{3}{30}+7\cdot\frac{3}{30}+8\cdot\frac{3}{30}
 =\frac{144}{30}=4.8.
\]

\item 甲厂产品的性价比为
\[
\frac{E(X_1)}{6}=\frac{6}{6}=1.
\]
乙厂产品的性价比为
\[
\frac{E(X_2)}{4}=\frac{4.8}{4}=1.2.
\]
因为 \(1.2>1\)，所以乙厂的产品更具可购买性.
\end{enumerate}
\end{solution}

\begin{problem}
如图，四棱锥 \(P-ABCD\) 中，\(PA\perp\) 底面 \(ABCD\)，四边形 \(ABCD\) 中，\(AB\perp AD\)，\(AB+AD=4\)，\(CD=\sqrt{2}\)，\(\angle CDA=45^\circ\).

\begin{enumerate}
\item 求证：平面 \(PAB\perp\) 平面 \(PAD\)；
\item 设 \(AB=AP\).

\begin{enumerate}
\item 若直线 \(PB\) 与平面 \(PCD\) 所成的角为 \(30^\circ\)，求线段 \(AB\) 的长；
\item 在线段 \(AD\) 上是否存在一个点 \(G\)，使得点 \(G\) 到点 \(P\)，\(B\)，\(C\)，\(D\) 的距离都相等？说明理由.
\end{enumerate}
\end{enumerate}

\bitmapfigure[max width=0.42\linewidth]{img/2011/fujian_science/q20_fig1.png}
\end{problem}

\begin{solution}
\begin{enumerate}
\item 在平面 \(PAD\) 内，直线 \(PA\) 与 \(AD\) 相交. 又因为 \(PA\perp\) 底面 \(ABCD\)，所以 \(PA\perp AB\)；题设还给出 \(AB\perp AD\). 因此直线 \(AB\) 垂直于平面 \(PAD\)，从而平面 \(PAB\perp\) 平面 \(PAD\).

\item 设 \(AB=u\)，\(AD=v\). 由题设，\(u+v=4\)，且 \(AP=AB=u\). 建立空间直角坐标系，取 \(A=(0,0,0)\)，\(B=(u,0,0)\)，\(D=(0,v,0)\)，\(P=(0,0,u)\). 由 \(CD=\sqrt{2}\)、\(\angle CDA=45^\circ\) 及图示位置，取 \(\overrightarrow{DC}=(1,-1,0)\)：它的长度为 \(\sqrt{2}\)，且与 \(\overrightarrow{DA}=(0,-v,0)\) 的夹角为 \(45^\circ\). 因此 \(C=(1,v-1,0)\).

\begin{enumerate}
\item 在平面 \(PCD\) 内取两条相交向量 \(\overrightarrow{PD}=(0,v,-u)\)，\(\overrightarrow{PC}=(1,v-1,-u)\). 于是平面 \(PCD\) 的一个法向量为
\[
\bs{n}=\overrightarrow{PD}\times\overrightarrow{PC}=(-u,-u,-v).
\]
直线 \(PB\) 的方向向量可取 \(\bs{d}=(1,0,-1)\). 若直线 \(PB\) 与平面 \(PCD\) 所成的角为 \(\theta\)，则
\[
\sin\theta=\frac{|\bs{d}\cdot\bs{n}|}{|\bs{d}|\,|\bs{n}|}=\frac{|v-u|}{\sqrt{2}\sqrt{2u^2+v^2}}.
\]
令 \(\theta=30^\circ\)，得
\[
\frac{(v-u)^2}{2(2u^2+v^2)}=\frac{1}{4},
\]
即 \(2(v-u)^2=2u^2+v^2\)，化简得 \(v(v-4u)=0\). 因为 \(v>0\)，所以 \(v=4u\). 结合 \(u+v=4\)，得
\[
u=\frac{4}{5},\qquad v=\frac{16}{5}.
\]
故 \(AB=\frac{4}{5}\).

\item 设 \(G=(0,t,0)\) 在线段 \(AD\) 上，则 \(0\leqslant t\leqslant v\). 若 \(G\) 到 \(P\)，\(B\)，\(C\)，\(D\) 的距离都相等，首先必须有 \(GP=GD\). 由 \(GP^2=t^2+u^2\)，\(GD^2=(v-t)^2\)，得
\[
t=\frac{v^2-u^2}{2v}=\frac{(16/5)^2-(4/5)^2}{2\cdot16/5}=\frac{3}{2}.
\]
此时 \(GB^2=t^2+u^2=GP^2\)，但
\[
GP^2=\left(\frac{3}{2}\right)^2+\left(\frac{4}{5}\right)^2=\frac{289}{100},
\]
而
\[
GC^2=1+\left(v-1-t\right)^2=1+\left(\frac{16}{5}-1-\frac{3}{2}\right)^2=\frac{149}{100}.
\]
所以 \(GC\ne GP\). 满足 \(GP=GD\) 的点已经唯一，故在线段 \(AD\) 上不存在满足条件的点 \(G\).
\end{enumerate}
\end{enumerate}
\end{solution}

\section{选考题：解答题}

\begin{problem}
设矩阵 \(M = \begin{pmatrix} a & 0 \\ 0 & b \end{pmatrix}\)（其中 \(a > 0\)，\(b > 0\)）

\begin{enumerate}
\item 若 \(a = 2\)，\(b = 3\)，求矩阵 \(M\) 的逆矩阵 \(M^{-1}\)；
\item 若曲线 \(C: x^{2} + y^{2} = 1\) 在矩阵 \(M\) 所对应的线性变换作用下得到曲线
\(C' : \frac{x^{2}}{4} + y^{2} = 1\)，求 \(a\)，\(b\) 的值.
\end{enumerate}
\end{problem}

\begin{solution}
\begin{enumerate}
\item 当 \(a=2\)，\(b=3\) 时，
\[
M=\begin{pmatrix}2&0\\0&3\end{pmatrix}.
\]
设其逆矩阵为 \(M^{-1}=\begin{pmatrix}r&0\\0&s\end{pmatrix}\)，由 \(MM^{-1}=I\) 得 \(2r=1\)，\(3s=1\). 所以
\[
M^{-1}=\begin{pmatrix}\frac{1}{2}&0\\0&\frac{1}{3}\end{pmatrix}.
\]

\item 矩阵 \(M\) 所对应的线性变换把原点保持不变，并把点 \((x_0,y_0)\) 变为 \((ax_0,by_0)\). 设变换后的点坐标为 \((x,y)\)，则 \(x_0=\frac{x}{a}\)，\(y_0=\frac{y}{b}\). 代入原曲线 \(x_0^2+y_0^2=1\)，得到变换后的曲线方程
\[
\frac{x^2}{a^2}+\frac{y^2}{b^2}=1.
\]
与题设曲线 \(\frac{x^2}{4}+y^2=1\) 比较，得 \(a^2=4\)，\(b^2=1\). 由于 \(a>0\)，\(b>0\)，所以 \(a=2\)，\(b=1\).
\end{enumerate}
\end{solution}

\begin{problem}
在直角坐标系 \(xOy\) 中，直线 \(l\) 的方程为 \(x-y+4=0\)，曲线 \(C\) 的参数方程为 \(\begin{cases}
x=\sqrt{3}\cos\alpha,\\
y=\sin\alpha
\end{cases}\)（\(\alpha\) 为参数）.

\begin{enumerate}
\item 已知在极坐标（与直角坐标系 \(xOy\) 取相同的长度单位，且以原点 \(O\) 为极点，以 \(x\) 轴正半轴为极轴）中，点 \(P\) 的极坐标为 \(\left(4,\frac{\pi}{2}\right)\)，判断点 \(P\) 与直线 \(l\) 的位置关系；
\item 设点 \(Q\) 是曲线 \(C\) 上的一个动点，求它到直线 \(l\) 的距离的最小值.
\end{enumerate}
\end{problem}

\begin{solution}
\begin{enumerate}
\item 极坐标 \(\left(4,\frac{\pi}{2}\right)\) 对应的直角坐标为
\[
P=\left(4\cos\frac{\pi}{2},4\sin\frac{\pi}{2}\right)=(0,4).
\]
代入直线 \(l\) 的方程，得 \(0-4+4=0\)，所以点 \(P\) 在直线 \(l\) 上.

\item 曲线 \(C\) 上的点 \(Q\) 可表示为 \(Q=(\sqrt{3}\cos\alpha,\sin\alpha)\). 点 \(Q\) 到直线 \(l\) 的距离为
\[
d=\frac{|\sqrt{3}\cos\alpha-\sin\alpha+4|}{\sqrt{2}}.
\]
由辅助角公式，
\[
\sqrt{3}\cos\alpha-\sin\alpha=2\cos\left(\alpha+\frac{\pi}{6}\right)
\]
所以其取值范围为 \([-2,2]\). 因此
\[
2\leqslant\sqrt{3}\cos\alpha-\sin\alpha+4\leqslant6
\]
绝对值可去掉，故
\[
d\geqslant\frac{2}{\sqrt{2}}=\sqrt{2}.
\]
当 \(\cos\left(\alpha+\frac{\pi}{6}\right)=-1\) 时取等号，所以距离的最小值为 \(\sqrt{2}\).
\end{enumerate}
\end{solution}

\begin{problem}
设不等式 \(|2x-1|<1\) 的解集为 \(M\).

\begin{enumerate}
\item 求集合 \(M\)；
\item 若 \(a\)，\(b\in M\)，试比较 \(ab+1\) 与 \(a+b\) 的大小.
\end{enumerate}
\end{problem}

\begin{solution}
\begin{enumerate}
\item 由绝对值不等式，
\[
-1<2x-1<1.
\]
不等式两边同时加 \(1\)，再除以正数 \(2\)，得 \(0<x<1\). 所以 \(M=(0,1)\).

\item 因为 \(a\)，\(b\in M\)，所以 \(0<a<1\)，\(0<b<1\). 于是
\[
ab+1-(a+b)=(1-a)(1-b)>0.
\]
因此 \(ab+1>a+b\).
\end{enumerate}
\end{solution}
